\section{A deterministic error model}
\label{sec:errors}

A bus that always delivers every byte in the right order is not a bus, but a function call with
extra steps. Real buses drop bytes, flip bits and delay deliveries.

The problem with testing error handling is that real errors are random. A test that fails
sometimes says nothing, and an error that happened once cannot be reproduced.

\textbf{The solution is to make the errors deterministic.} The bus implementation in the course's
simulation counts bytes and does exactly the same thing on every run:

\begin{cppcode}
// Deterministic fault injection: byte five is dropped, every run, every time.
if (5U == ++myByteCounter) { return; }
\end{cppcode}

Three errors are enough to cover most of it:

\begin{booktable}{@{}lL@{}}
\thead{Error} & \thead{What happens} \\ \midrule
\textbf{Dropped byte} & Every following field is shifted. The parser reads the wrong length,
  waits for the wrong number of bytes, and finally fails on the checksum. \\
\textbf{Flipped bit} & The frame is structurally correct --- right length, fields in the right
  places --- but the checksum does not add up. That is precisely the case the checksum exists
  for. \\
\textbf{Delay} & The frame arrives, whole and correct, but later than expected. Nothing breaks,
  and that is the point: a correct parser does not care \emph{when} the bytes arrive. \\
\end{booktable}

The third case is worth pausing at. A parser that works when the bytes arrive close together but
breaks when they arrive far apart has a hidden timing dependency, and that is a bug which never
shows itself until the system is moved to slower hardware.
